% Auto-generated by Code2Latex
% Requires: longtable, booktabs packages

\begin{longtable}{p{\textwidth}}
\caption{Tools for wetlab (Level 1)} \\
\endfirsthead

\multicolumn{1}{c}{\tablename\ \thetable{} -- continued from previous page} \\
\endhead

\multicolumn{1}{r}{Continued on next page} \\
\endfoot

\endlastfoot

\textbf{\texttt{add\_a\_solution}} \\
\begin{description}
    \item[\textbf{Description:}] Adds a specific volume of `sol2\_label` to all of `sol1\_label` and returns observations about the changes of solution color and precipitation amount and color.

    \item[\textbf{Arguments:}]
    \begin{description}
        \item \texttt{test\_label} (str): label of the resulting solution
        \item \texttt{sol1\_label} (str): label of the first solution
        \item \texttt{sol2\_label} (str): label of the second solution
        \item \texttt{sol2\_vol} (int): volume of the second solution to draw
    \end{description}
    \item[\textbf{Returns:}] a string containing the observations from the test
\end{description} \\
\midrule
\textbf{\texttt{add\_precipitate\_to\_solution}} \\
\begin{description}
    \item[\textbf{Description:}] Adds a precipitate to a specific volume of a solution and returns observations about the changes in the amount/color of the added precipitate or the solution color.

    \item[\textbf{Arguments:}]
    \begin{description}
        \item \texttt{test\_label} (str): label of the resulting mixture
        \item \texttt{prec\_label} (str): label of the precipitate
        \item \texttt{sol\_label} (str): label of the solution
        \item \texttt{sol\_vol} (int): volume of the solution
    \end{description}
    \item[\textbf{Returns:}] a string containing the observations from the test
\end{description} \\
\midrule
\textbf{\texttt{check\_inventory}} \\
\begin{description}
    \item[\textbf{Description:}] Returns the current contents of the Inventory
    \item[\textbf{Arguments:}]
    None
    \item[\textbf{Returns:}] a string containing the solutions and precipitates in the Inventory
\end{description} \\
\midrule
\textbf{\texttt{checkout\_color}} \\
\begin{description}
    \item[\textbf{Description:}] Observe the color of a solution or precipitate.
    \item[\textbf{Arguments:}]
    \begin{description}
        \item \texttt{label} (str): label of the target object
    \end{description}
    \item[\textbf{Returns:}] the type (reagent, solution, or precipitate) and the color of the object
\end{description} \\
\midrule
\textbf{\texttt{filter\_solution}} \\
\begin{description}
    \item[\textbf{Description:}] Separates the precipitate from the supernatant solution.
    \item[\textbf{Arguments:}]
    \begin{description}
        \item \texttt{label} (str): label of the target solution
    \end{description}
    \item[\textbf{Returns:}] a string with a message about the success/failure of the filtration
\end{description} \\
\midrule
\textbf{\texttt{get\_available\_reagents}} \\
\begin{description}
    \item[\textbf{Description:}] Returns the list of available reagent solutions.
    \item[\textbf{Arguments:}]
    None
    \item[\textbf{Returns:}] a string containing all available reagents
\end{description} \\
\midrule
\textbf{\texttt{lookup\_flame\_colors}} \\
\begin{description}
    \item[\textbf{Description:}] Returns the characteristic flame colors of cations.
    \item[\textbf{Arguments:}]
    None
    \item[\textbf{Returns:}] the list of characteristic flame colors
\end{description} \\
\midrule
\textbf{\texttt{lookup\_precipitate\_colors}} \\
\begin{description}
    \item[\textbf{Description:}] Returns the list of colored precipitates and their color.
    \item[\textbf{Arguments:}]
    None
    \item[\textbf{Returns:}] the list of precipitate colors
\end{description} \\
\midrule
\textbf{\texttt{measure\_pH}} \\
\begin{description}
    \item[\textbf{Description:}] Measures the pH of the solution using a pH paper.
    \item[\textbf{Arguments:}]
    \begin{description}
        \item \texttt{label} (str): label of the target solution
    \end{description}
    \item[\textbf{Returns:}] the closest integer value to the actual pH of the solution
\end{description} \\
\midrule
\textbf{\texttt{mix\_two\_solutions}} \\
\begin{description}
    \item[\textbf{Description:}] Mixes two solutions with the given volumes an returns observations about precipitation and color of the resulting solution
    \item[\textbf{Arguments:}]
    \begin{description}
        \item \texttt{test\_label} (str): label of the resulting solution
        \item \texttt{sol1\_label} (str): label of the first solution
        \item \texttt{sol1\_vol} (int): volume of the first solution to draw
        \item \texttt{sol2\_label} (str): label of the second solution
        \item \texttt{sol2\_vol} (int): volume of the second solution to draw
    \end{description}
    \item[\textbf{Returns:}] a string containing the observations from the test
\end{description} \\
\midrule
\textbf{\texttt{perform\_flame\_test}} \\
\begin{description}
    \item[\textbf{Description:}] Performs a flame test on the solution. Cost = 1 mL
    \item[\textbf{Arguments:}]
    \begin{description}
        \item \texttt{label} (str): label of the target solution
    \end{description}
    \item[\textbf{Returns:}] the resulting observation from the flame test
\end{description} \\
\midrule
\textbf{\texttt{possible\_anions}} \\
\begin{description}
    \item[\textbf{Description:}] Returns the list of possible anions.
    \item[\textbf{Arguments:}]
    None
    \item[\textbf{Returns:}] a string containing all the possible anions
\end{description} \\
\midrule
\textbf{\texttt{possible\_cations}} \\
\begin{description}
    \item[\textbf{Description:}] Returns the list of possible cations.
    \item[\textbf{Arguments:}]
    None
    \item[\textbf{Returns:}] a string containing all the possible cations
\end{description} \\
\midrule
\textbf{\texttt{simulate\_color\_mixture}} \\
\begin{description}
    \item[\textbf{Description:}] Given a mixture of precipitate colors, it mixes them with the given fractions and returns the name of the resulting precipitate color.
    \item[\textbf{Arguments:}]
    \begin{description}
        \item \texttt{mixture} (list[tuple[str, float]]): mixture components and their fractions as a list of tuples
    \end{description}
    \item[\textbf{Returns:}] name of the resulting color
\end{description} \\
\end{longtable}
